% =====================================================================
% Subsection to close Section III (The SHIFT-SURV Benchmark).
% Place immediately BEFORE "A no-shift control uses the same parameter
% tuple for source and target centers." — that sentence stays where it
% is, at the end of III-C, and must not be repeated here.
%
% Revision notes (prof's corrections, all incorporated):
%   1. Baseline-hazard shift preserves the LATENT ordering but can still
%      move the numerical AUROC at fixed tau, because it changes the
%      composition of the event / non-event groups. Verified: ERM target
%      minus source AUROC is +0.018 under baseline hazard against -0.011
%      for the no-shift control, and +0.003 under label prior. Small and
%      not individually significant at ten seeds, hence the modal "can".
%   2. Measurement shift does NOT change the latent ordering — only the
%      map from latent factors to covariates, hence the ordering
%      available to a predictor trained on X. Collapsing it with concept
%      shift would destroy the recoverable / unrecoverable distinction
%      the two-reference decomposition rests on.
%   3. All six regimes are now covered, grouped by mechanism rather than
%      listed by name. Grouping order is deliberately NOT Table I order:
%      it escalates in severity and ends where the two-reference
%      decomposition begins.
%
% Verified against the generator:
%   - covariate: _covariate_shift uses _orthogonal_to(w_src), so
%     w^T mu_inv_tgt = 0 and the eta distribution (hence prevalence) is
%     unchanged.  scenarios.py:277
%   - spurious: the diagonal blocks of latent_cov stay I, so the
%     marginal of u_inv is identical in source and target; only the
%     cross-covariance C_d differs.  scenarios.py:393-413
%   - measurement: _observational_shift inherits mu and sigma (defaults
%     are 0.0 / 0.0 / None), so R is a pure orthogonal rotation with no
%     information destroyed.  scenarios.py:~445
%
% Space note: the label-prior and concept display equations are already
% in III-C and are cross-referenced rather than repeated. If the page
% budget bites, drop \begin{itemize}/\end{itemize} and run the four
% blocks as consecutive sentences — recovers ~0.35 column.
% =====================================================================

\subsection{Discrimination and calibration}
\label{subsec:disc-cal}

A transported risk model has two distinct targets: preserving the ordering of
patients by risk, and preserving the probabilistic meaning of its predictions.
SHIFT-SURV perturbs these targets through distinct controlled mechanisms, so
they are separable at the level of the data-generating process, although the
empirical measures of the two need not be statistically independent.

\emph{Discrimination} is the ability to rank patients by risk. For a predictor
$f_\theta$ evaluated on the target center,
\begin{equation}
\mathrm{AUROC}(f_\theta)
=
P\!\left(
f_\theta(X_i) > f_\theta(X_j)
\;\middle|\;
Y_i^{(\tau)}=1,\,
Y_j^{(\tau)}=0
\right),
\end{equation}
with the usual half-weight assigned to ties. AUROC depends on the predictor
only through the ordering it induces, and is therefore unchanged by strictly
increasing transformations of its scores.

\emph{Calibration} is the agreement between predicted probabilities and
observed event frequencies. A predictor is calibrated on the target when
$P(Y^{(\tau)}=1\mid f_\theta(X))=f_\theta(X)$ almost surely. We summarize
departures from this condition using the expected calibration error over $B$
equal-width probability bins,
\begin{equation}
\mathrm{ECE}(f_\theta)
=
\sum_{b=1}^{B}
\frac{n_b}{n}
\left|\bar{y}_b-\bar{q}_b\right|,
\end{equation}
where $n_b$ is the number of target patients whose predicted probability falls
in bin $b$, and $\bar{y}_b$ and $\bar{q}_b$ are the observed event rate and the
mean predicted probability in that bin. The two properties are complementary
rather than redundant: a model can preserve patient ranking while assigning
systematically miscalibrated probabilities, and predictions can carry the
correct average event rate while providing little ranking information.
Discrimination is what prioritization requires, whereas decisions taken at a
threshold on absolute risk also require calibration.

The generator makes the distinction explicit. The fixed-horizon event
probability is
\begin{equation}
P\!\left(Y_{i,d}^{(\tau)}=1 \mid u_{i,d}^{\mathrm{inv}}\right)
=
1-\exp\!\left\{
-\lambda_{0,d}\,\tau^{\nu_d}\exp(\eta_{i,d})
\right\},
\end{equation}
so the coefficient vector $w_d$ fixes the latent ordering through
$\eta_{i,d}$, the observation map $A_d^{\mathrm{inv}}$ fixes how much of that
ordering is recoverable from $X$, and $\lambda_{0,d}$ and $\nu_d$ set the level
at which the ordering is converted into probabilities.

The regimes separate accordingly.
\begin{itemize}
\item \emph{Covariate shift} moves the latent means, but along a direction
orthogonal to $w_d$, so the distribution of $\eta_{i,d}$ and the event
prevalence at $\tau$ are unchanged. Neither the map from covariates to risk nor
the level at which it is converted into probabilities is affected; only the
region of covariate space that is populated moves.

\item \emph{Baseline-hazard} and \emph{label-prior shift} preserve the latent
ordering and change the level. A change in $\lambda_{0,d}$ alters the target
event rate, and hence the calibration of a source-trained predictor; it can
nevertheless alter the numerical AUROC at a fixed horizon, because a different
baseline hazard changes the composition of the event and non-event groups at
$\tau$. Label-prior shift preserves the class-conditional covariate
distributions by construction, so the population AUROC of a fixed scoring rule
is preserved exactly while its calibration changes.

\item \emph{Measurement} and \emph{spurious-association shift} preserve the
latent ordering but not the ordering a source-trained predictor recovers from
$X$. Measurement shift leaves $\eta_{i,d}$ untouched and rotates the map from
latent factors to covariates. Spurious shift leaves both $\eta_{i,d}$ and the
observation map untouched, so the target-optimal rule on $X$ is unchanged; what
differs is the rule a source-trained model estimates, since the spurious block
of $X$ proxies for the invariant signal in the source and not in the target. In
both regimes the target-optimal rule remains a function of $X$, and the gap is
therefore recoverable in principle.

\item \emph{Concept shift} changes $w_d$ and therefore the latent ordering
itself. The target relationship is learnable from labeled target data, but is
not identified by any source-trained rule, whatever the observation process.
\end{itemize}
The last two groups are what the two reference models of
Section~\ref{sec:setup} are designed to separate, and both references are
evaluated on discrimination and calibration alike, so the decomposition applies
to either axis.
